\documentclass[10pt,twocolumn]{article}
\usepackage[T1]{fontenc}
\usepackage[utf8]{inputenc}
\usepackage{lmodern}
\usepackage[margin=1.8cm,top=2.4cm,bottom=2cm,columnsep=0.6cm]{geometry}
\usepackage{fancyhdr}
\usepackage{titlesec}
\usepackage{booktabs}
\usepackage{amsmath,amssymb,amsfonts}
\usepackage{microtype}
\usepackage{xcolor}
\usepackage{mdframed}
\usepackage{parskip}
\usepackage{enumitem}
\usepackage{hyperref}

\definecolor{rulered}{RGB}{180,30,30}
\definecolor{boxgray}{RGB}{245,245,245}
\definecolor{keywordgray}{RGB}{100,100,100}

\pagestyle{fancy}
\fancyhf{}
\fancyhead[L]{\textsc{\small Claude Mag}}
\fancyhead[C]{\small\textit{Machine Learning Research Digest}}
\fancyhead[R]{\small 2026-10-21}
\fancyfoot[C]{\small\thepage}
\renewcommand{\headrulewidth}{0.8pt}

\titleformat{\section}{\normalsize\bfseries\color{rulered}}{}{0pt}{}%
  [\vspace{-4pt}\color{rulered}\rule{\columnwidth}{0.4pt}]
\titlespacing{\section}{0pt}{8pt}{4pt}

\hypersetup{colorlinks=true,urlcolor=rulered,linkcolor=black}

\begin{document}
\setlength{\parindent}{0pt}
\setlength{\parskip}{4pt}

{\small\color{keywordgray}\textbf{Keywords:}
  multi-agent systems \textbullet{}
  LLM agents \textbullet{}
  goal drift \textbullet{}
  reinforcement learning}\par\vspace{4pt}

{\LARGE\bfseries\raggedright
  Learning to Steer LLM Agents and Multi-Agent Systems}\par\vspace{4pt}

{\large\itshape\raggedright
  A Dual-Path Steering Agent Trained with GRPO Corrects Goal Drift
  and Transfers to Unseen Executors}\par\vspace{6pt}

{\small\textbf{By} Leyang Shen, Yaxin Gong, et al.}\par\vspace{2pt}

{\small\color{keywordgray}arXiv:2610.22316 \textbullet{}
  October 2026}
\par\vspace{2pt}\noindent{\color{rulered}\rule{\columnwidth}{0.6pt}}\vspace{6pt}

LLM-based agents—whether acting alone or in concert—are prone to goal
drift: as execution unfolds step by step, the system progressively deviates
from its original objective. This paper introduces a dedicated \textit{steering
agent} that monitors ongoing execution and corrects drift via two intervention
paths, jointly trained with Group Relative Policy Optimization (GRPO). On
information-seeking benchmarks, the steered system completes tasks more
accurately and efficiently than executor-only baselines, and its policy
transfers to unseen task types and executor architectures without retraining.

\begin{mdframed}[backgroundcolor=boxgray,linecolor=rulered,linewidth=1pt,
    innerleftmargin=6pt,innerrightmargin=6pt,
    innertopmargin=6pt,innerbottommargin=6pt]
\textbf{\small AT A GLANCE}\par\vspace{4pt}
\begin{description}[leftmargin=0pt,labelsep=4pt,itemsep=2pt,parsep=0pt,
    font=\small\bfseries]
  \item[Problem:] Agents and multi-agent pipelines drift from the overall
    goal as execution steps accumulate, wasting compute and producing
    unintended outcomes.
  \item[Why it matters:] No existing component has global oversight of
    execution trajectory; prompt-based corrections lack a trained policy
    for when and how to intervene.
  \item[Method:] Dedicated steering agent with fast (cheap) and slow
    (expensive) intervention paths, jointly trained with GRPO and a
    pairwise ranking objective.
  \item[Result:] Improved task completion rate and fewer total execution
    steps; policy transfers to out-of-domain tasks and unseen executors.
  \item[Evidence:] Experiments on information-seeking single-agent and
    multi-agent tasks; ablations confirm both paths are necessary.
\end{description}
\end{mdframed}

\section{The Big Picture}

Modern LLM agent pipelines are long: they involve dozens of tool calls,
web searches, sub-task handoffs, and intermediate reasoning steps. Each
step introduces the possibility that the agent drifts from the high-level
objective—optimizing a sub-task at the expense of the whole, misinterpreting
a retrieved document, or losing track of a constraint stated many steps
earlier.

In multi-agent settings the problem compounds: agents optimize locally,
communication between them is asynchronous, and no single agent has a
global view of the pipeline's progress.

This paper's central argument is that the solution should not be to redesign
every executor agent but to add a separate \textit{steering} agent whose
sole job is global oversight. The steering agent watches, decides when to
act, and intervenes minimally to keep the execution on course.

\section{Why Existing Approaches Fall Short}

\textbf{Prompt injection:} Injecting correction instructions into the next
step is brittle because it lacks a policy for deciding when or how to
intervene. Over-intervention stalls progress; under-intervention allows drift
to persist.

\textbf{Reward shaping during training:} Re-training all executors with a
global objective requires access to each agent's parameters and is
incompatible with off-the-shelf or proprietary executors.

\textbf{Post-hoc replanning:} Replanning after detecting failure is costly:
it requires abandoning completed work and restarting. It cannot prevent drift
mid-execution.

\section{Core Mechanism}

The steering agent operates alongside the executor(s) without modifying them.
At each execution step it observes the full context and chooses among three
options:

\begin{enumerate}[itemsep=2pt,parsep=0pt]
  \item \textbf{Silent:} No intervention; execution continues.
  \item \textbf{Fast path (low-cost):} A lightweight corrective signal,
    such as re-anchoring the executor to the current sub-goal or pruning
    off-track tool calls.
  \item \textbf{Slow path (high-cost):} A full re-planning step or a
    coordinator message to all agents, for novel or complex deviations.
\end{enumerate}

Both paths and the routing decision are trained jointly, so the steering
agent learns not only how to correct but also when each type of correction
is warranted.

\section{Technical Formulation}

Let $s_t$ be the execution state at step $t$ and $a^{\text{exec}}_t$ the
action the executor would take absent intervention. The steering agent
produces $\delta_t \in \{\text{silent}, \text{fast}, \text{slow}\}$ and
a steering signal $\sigma_t$. The executor's effective action is:
\[
a_t = f(a^{\text{exec}}_t,\; \sigma_t)
\]
where $f$ merges the executor's action with the steering signal. The GRPO
objective for the steering agent is:
\[
\mathcal{J}_{\text{steer}}(\theta) = \mathbb{E}\!\left[\sum_t \!\left(R_{\text{final}}
  - \lambda_{\text{cost}} \cdot \mathbf{1}[\delta_t = \text{slow}]\right)
  \cdot \log \pi_\theta(\delta_t, \sigma_t \mid s_t)\right]
\]
with $\lambda_{\text{cost}}$ penalising slow-path invocations, pushing the
policy to use the cheap fast path whenever sufficient.

A supplementary pairwise ranking loss compares trajectory pairs (same task,
different steering decisions) and penalises whichever produced the worse
outcome:
\[
\mathcal{L}_{\text{rank}}(\theta) = \mathbb{E}_{(\tau^+, \tau^-)}\!\left[
  \max\!\left(0,\; 1 - \log\frac{\pi_\theta(\tau^+)}{\pi_\theta(\tau^-)}\right)\right]
\]
The final training objective is $\mathcal{J}_{\text{steer}} - \alpha \mathcal{L}_{\text{rank}}$.

\section{Learning or Inference Procedure}

\begin{enumerate}[itemsep=2pt,parsep=0pt]
  \item Collect paired trajectories: one with the steering agent active,
    one without (or with a different steering decision).
  \item Annotate outcomes; form pairwise preference labels for the ranking
    loss.
  \item Train the steering agent with the combined GRPO + ranking objective.
  \item At inference, the steering agent runs alongside the executor(s);
    no modification to executors is required.
  \item Intervention frequency is controlled post-training via
    $\lambda_{\text{cost}}$, enabling efficiency-accuracy trade-offs.
\end{enumerate}

\section{What the Guarantee Says}

The pairwise ranking objective provides a contrastive guarantee: the steering
policy converges to prefer interventions whose outcomes are better in expectation
than the counterfactual of no intervention. The GRPO component constrains the
policy to stay close to the reference model, preventing exploitation of the
steering interface. Together, the objectives bound the expected cost overhead
of incorrect interventions relative to executor-only execution.

\section{Experimental Findings}

\begin{itemize}[itemsep=2pt,parsep=0pt]
  \item \textbf{Single-agent information-seeking:} The steering agent
    improves task completion rate and reduces total execution steps
    compared to the executor-only baseline, prompt-injection correction,
    and replanning-on-failure.
  \item \textbf{Multi-agent pipelines:} Gains are larger in multi-agent
    settings, where drift compounds across agents; the steering agent
    provides global coordination that individual agents lack.
  \item \textbf{Transfer to out-of-domain tasks:} The steering policy,
    trained on information-seeking tasks, transfers to document QA and
    planning tasks without any fine-tuning.
  \item \textbf{Transfer to unseen executors:} Executors not encountered
    during training are steered effectively, confirming that the policy
    learned general drift-correction rather than executor-specific heuristics.
\end{itemize}

\section{Ablations and Interpretation}

\textbf{Fast-path only:} Loses on novel drift patterns requiring deep
re-planning; confirms the slow path is necessary for non-routine corrections.

\textbf{Slow-path only:} Higher cost per step with marginal accuracy gain
over fast-only; the fast path handles the majority of cases cheaply.

\textbf{Without pairwise ranking (GRPO alone):} Noisier intervention
decisions; the steering agent more often invokes the slow path when the fast
path would have sufficed, indicating the ranking loss sharpens the routing
decision.

\textbf{Controllable intervention:} Sweeping $\lambda_{\text{cost}}$ traces
a smooth Pareto curve between task accuracy and intervention cost, confirming
the design supports post-training trade-off adjustment.

\section*{Reference}

Leyang Shen, Yaxin Gong, et al.
\textbf{Learning to Steer LLM Agents and Multi-Agent Systems}.
arXiv:2610.22316, October 2026.
\url{https://arxiv.org/abs/2610.22316}

\end{document}
